% Part: first-order-logic
% Chapter: beyond
% Section: other-logics

\documentclass[../../../include/open-logic-section]{subfiles}

\begin{document}

\olfileid{fol}{byd}{oth}

\olsection{Logika Liyane}

Kaya kang wis bisa sampeyan kira,
ngrancang logika anyar ora angel.
Sampeyan dhewe uga bisa nggawe
sintaks, nyusun sistem deduktif,
lan ngrancang semantik kang cocog.
Yen sampeyan kepengin sistem
!!{derivation} iku jangkep kanggo
semantike, mbokmenawa sampeyan
kudu rada pinter golek cara.
Sampeyan uga kudu ngupaya
supaya wong akeh yakin yen
logika sampeyan pancen narik
kawigaten. Nanging minangka
ganjarane, sampeyan bisa
ngentekake akeh wektu kanthi
seneng-seneng nliti sipat
matematis lan komputasional
logika kasebut.

Ing dasawarsa-dasawarsa pungkasan,
logika formal saya akeh lan
warna-warna. Logika samar
dirancang kanggo nggambarake
nalar ngenani sipat kang ora
temenan cetha. Logika
probabilistik dirancang
kanggo nggambarake nalar
ngenani bab kang durung mesthi.
Logika default lan logika
nonmonoton dirancang kanggo
nggambarake nalar kang
kesimpulane bisa dibatalake,
yaiku inferensi kang "lumrah"
nanging bisa digugurake maneh
nalika ana katrangan anyar.
Ana logika epistemik kanggo
nggambarake nalar ngenani
kawruh; logika kausal
kanggo nalar ngenani
sesambungan sebab-akibat;
lan malah logika "deontik"
kanggo nalar ngenani
kewajiban moral lan etika.
Gumantung saka apa sebab
utama sistem-sistem iki
dikenalake---filsafat,
matematika, utawa
komputasi---sampeyan bisa
nemokake panliten ngenani
logika kasebut ing
logika matematika,
logika filsafat,
kacerdhasan gawean,
ilmu kognitif,
utawa papan liyane.

Dhaptar iki isih bisa
diterusake, lan kamungkinane
kaya ora ana enteke.
Mbokmenawa kita ora bakal
kasil nggayuh impen Leibniz
ngowahi kabeh nalar manungsa
dadi petungan---nanging
iku ora bisa ngalang-alangi
kita kanggo nyoba.

\end{document}
